% 三条历史访问路径；最终宽度约162mm，按保存粒度比较。
\begin{tikzpicture}[
  font=\figurefont\footnotesize,
  rsbox/.style={rectangle,draw=BookRule,fill=BookPaper,
    text=BookInk,align=center,text width=40mm,minimum height=13mm,inner sep=4pt},
  rsflow/.style={-{Stealth[length=2mm]},draw=BookMuted,line width=.9pt}
]
  \node[font=\figurefont\footnotesize\bfseries] at (0,1.2) {新事件到来};
  \node[font=\figurefont\footnotesize\bfseries] at (5.7,1.2) {保留的信息};
  \node[font=\figurefont\footnotesize\bfseries] at (11.4,1.2) {预测请求读取};
  \node[rsbox,fill=FigureInput!10] (a1) at (0,0) {增量读写\\擦除、添加、状态更新};
  \node[rsbox,fill=FigureModel!10] (b1) at (5.7,0) {固定容量记忆\\部分事件细节被压缩};
  \node[rsbox,fill=FigureOutput!10] (c1) at (11.4,0) {读取记忆及归纳状态\\结合候选预测};
  \node[rsbox,fill=FigureInput!10] (a2) at (0,-1.9) {追加原始事件\\更新历史索引};
  \node[rsbox,fill=FigureModel!10] (b2) at (5.7,-1.9) {可检索事件记录\\保留逐条行为信息};
  \node[rsbox,fill=FigureOutput!10] (c2) at (11.4,-1.9) {候选检索相关事件\\再作精细聚合};
  \node[rsbox,fill=FigureInput!10] (a3) at (0,-3.8) {更新哈希桶或行为簇\\刷新聚合及版本};
  \node[rsbox,fill=FigureModel!10] (b3) at (5.7,-3.8) {桶或簇表示\\保留聚合信息与数量};
  \node[rsbox,fill=FigureOutput!10] (c3) at (11.4,-3.8) {候选访问桶或簇\\形成聚合兴趣};
  \foreach \r in {1,2,3} {
    \draw[rsflow] (a\r.east) -- (b\r.west);
    \draw[rsflow] (b\r.east) -- (c\r.west);
  }
\end{tikzpicture}
